% ============================================================================
\section{Algoritmo constructivo}\label{sec:constructivo}
% ============================================================================

\subsection{Idea}

Un algoritmo constructivo parte de la solución vacía y añade columnas una a una hasta que
todas las filas están cubiertas. Si en cada paso se toma la columna que \emph{en ese momento}
parece mejor, sin volver nunca atrás, el algoritmo se llama \textbf{voraz} (en inglés
\emph{greedy}). Los apuntes lo ilustran con el ``vecino más próximo'' del problema del
viajante; aquí la pregunta es qué significa ``la columna que parece mejor''.

\subsection{Tres criterios posibles}

Sea $g_j$ el número de filas \emph{todavía sin cubrir} que cubriría la columna $j$ (su
\emph{ganancia}). Hay tres criterios naturales:
\begin{center}
\begin{tabular}{@{}cll@{}}
\toprule
Criterio & Se elige la columna con\ldots & Inconveniente \\ \midrule
1 & menor coste $c_j$ & ignora cuánto cubre \\
2 & mayor ganancia $g_j$ & ignora cuánto cuesta \\
3 & mayor cociente $g_j/c_j$ & --- \\
\bottomrule
\end{tabular}
\end{center}
El criterio 3 elige la columna que cubre más filas nuevas \emph{por unidad de coste}; dicho
de otro modo, la que cubre cada fila nueva al menor precio. Es el voraz clásico del
cubrimiento (Chvátal, 1979) y tiene una garantía teórica: su solución nunca cuesta más de
$H_m = 1+\tfrac12+\dots+\tfrac1m \approx \ln m$ veces el óptimo~\cite[sec.~2.3]{kordalewski}.
En el vocabulario de los apuntes es, por tanto, un \emph{algoritmo aproximado}. En la
práctica queda muchísimo más cerca del óptimo que esa cota.

Un detalle esencial es que el criterio es \textbf{adaptativo}: $g_j$ cambia después de cada
inserción, porque las filas que se acaban de cubrir dejan de contar para las demás columnas.

\begin{algorithm}[H]
\caption{Constructivo voraz (criterio del cociente)}\label{alg:voraz}
\begin{algorithmic}[1]
\State $S\gets\emptyset$
\While{queda alguna fila sin cubrir}
  \State para cada columna $j\notin S$: $g_j\gets$ n.º de filas sin cubrir que cubre $j$
  \State $j^\ast\gets$ columna con $g_j>0$ que maximiza $g_j/c_j$
  \State $S\gets S\cup\{j^\ast\}$
\EndWhile
\State eliminar de $S$ las columnas redundantes, la más cara primero
\State \Return $S$
\end{algorithmic}
\end{algorithm}

\subsection{El paso final: quitar redundantes}

Una columna elegida al principio puede volverse innecesaria más tarde, cuando otras columnas
han cubierto todas sus filas. Por eso, al terminar, se eliminan las redundantes.

\begin{decision}[Decisión: quitar primero la redundante más cara]
Si hay varias redundantes no siempre se pueden quitar todas: dos columnas pueden ser
redundantes por separado y dejar de serlo en cuanto se quita la otra (si comparten una fila
que solo cubren ellas dos). Hay que decidir un orden. Se quita primero la más cara porque es
la que más ahorra, se recalcula, y se repite. Es, a su vez, un criterio voraz.
\end{decision}

\subsection{El algoritmo, a mano, sobre el ejemplo}

Con los datos de la sección~\ref{sec:ejemplo} ($c=(2,5,4,3,1)$) y el criterio 3:

\begin{center}
\small
\begin{tabular}{@{}clcccccl@{}}
\toprule
Paso & Filas sin cubrir & \multicolumn{5}{c}{Cociente $g_j/c_j$ de las columnas 1 a 5} & Se elige \\ \midrule
1 & 1, 2, 3 & $1/2=0.50$ & $2/5=0.40$ & $1/4=0.25$ & $2/3=0.67$ & $1/1=\mathbf{1.00}$ & columna 5 \\
2 & 1, 3 & $1/2=\mathbf{0.50}$ & $1/5=0.20$ & $1/4=0.25$ & $1/3=0.33$ & --- & columna 1 \\
3 & 3 & --- & $1/5=0.20$ & $1/4=\mathbf{0.25}$ & --- & --- & columna 3 \\
\bottomrule
\end{tabular}
\end{center}

Resultado: $S=\{5,1,3\}$, coste $1+2+4=7$, que es el óptimo. Nótese cómo cambian las
ganancias: en el paso 1 la columna 4 cubría dos filas nuevas; en el paso 2 solo una, porque
la fila 2 ya la había cubierto la columna 5. Con el criterio 2 (mayor ganancia) el mismo
algoritmo elige las columnas 4 y 2 y obtiene coste 8. Se puede comprobar con el programa:
\begin{verbatim}
   ./scp ejemplo.dat voraz criterio=3 ver=1
   ./scp ejemplo.dat voraz criterio=2 ver=1
\end{verbatim}

\subsection{Implementación: la función \cod{completar}}\label{sec:completar}

El constructivo es tan corto que cabe en tres líneas, porque todo el trabajo lo hace la
función común \cod{completar}:

\fragmento{constructivo.c}{24}{30}

\cod{completar(I, S, criterio, lrc, memoria, lambda)} cubre las filas que $S$ tenga sin
cubrir. Sus argumentos permiten usar la misma función en todos los algoritmos:
\begin{itemize}
  \item \cod{criterio}: 1, 2 o 3 (los de la tabla anterior);
  \item \cod{lrc}: tamaño de la \emph{lista restringida de candidatos}. En cada paso no se
        toma obligatoriamente la mejor columna, sino una al azar entre las \cod{lrc} mejores.
        Con \cod{lrc = 1} es el voraz puro, determinista. Con valores mayores es el voraz
        \emph{aleatorizado} que necesitan GRASP y las reparaciones (sección~\ref{sec:grasp});
  \item \cod{memoria} y \cod{lambda}: solo los usa GRASP.
\end{itemize}

El código tiene tres bloques. El primero calcula las ganancias iniciales recorriendo
únicamente las filas sin cubrir:

\fragmento{base.c}{493}{504}

El segundo, dentro del bucle principal, recorre las candidatas, descarta las que ya no
aportan nada y mantiene ordenadas las \cod{lrc} mejores (una inserción en una lista ordenada
muy corta):

\fragmento{base.c}{522}{549}

El tercero actualiza las ganancias \emph{antes} de insertar la columna elegida: por cada fila
que la elegida va a cubrir por primera vez, todas las columnas que cubren esa fila pierden
una unidad de ganancia.

\fragmento{base.c}{551}{560}

\begin{idea}[Por qué es rápido]
Las ganancias no se recalculan desde cero en cada paso: se \emph{actualizan}. Cada fila pasa
de ``sin cubrir'' a ``cubierta'' una sola vez, y en ese momento se recorre su lista de
columnas. En total se recorren los unos de la matriz como mucho una vez ($\approx 75\,000$
operaciones), más el repaso de la lista de candidatas en cada paso. Una construcción
completa tarda unos 8 milisegundos.
\end{idea}

\subsection{Comparación de los tres criterios}\label{sec:res-voraz}

La tabla~\ref{tab:basicos} (página~\pageref{tab:basicos}) recoge el resultado de los tres
criterios sobre las 30 instancias. La medida es la desviación porcentual respecto del óptimo
(RPD, definida en la sección~\ref{sec:rpd}). En las 6 instancias de calibración:
\begin{itemize}
  \item el criterio 1 (la más barata) queda un 15.9\,\% por encima del óptimo y usa unas 243
        columnas: llena la solución de columnas baratas que cubren muy poco;
  \item el criterio 2 (la que más cubre) usa solo 49 columnas pero queda un 33.4\,\% por
        encima: elige columnas grandes sin mirar que son carísimas;
  \item el criterio 3 (cociente) queda un 6.6\,\% por encima, con unas 167 columnas.
\end{itemize}

\begin{decision}[Decisión: criterio del cociente]
Se adopta el criterio 3 como constructivo de la práctica y como base de todas las
reparaciones y construcciones de las metaheurísticas. La elección se ha hecho mirando solo
las 6 instancias de calibración; las 24 de validación la confirman (6.4\,\% frente a
15.5\,\% y 31.6\,\%). Coincide con lo esperado por la estructura de las instancias (los
costes están correlacionados con la cobertura, sección~\ref{sec:instancias}).
\end{decision}
